% TOP_PROBLEM: {"problemId": "problem.segreharbournegimiglianohirschowitz-conjecture", "canonicalTitle": "Segre–Harbourne–Gimigliano–Hirschowitz Conjecture", "releaseRank": 493, "primaryDomain": "geometry_topology", "primaryDomainLabel": "Geometry & topology", "displayStatus": "open_with_solved_subcases", "releaseStatus": "open_with_solved_subcases"}
% !TeX program = lualatex
% Definition and sources only; this file is not an LLM solution attempt.
\documentclass[11pt]{article}
\usepackage[margin=25mm]{geometry}
\usepackage{fontspec}
\setmainfont{Latin Modern Roman}
\usepackage{amsmath,amssymb,mathtools}
\usepackage{unicode-math}
\setmathfont{Latin Modern Math}
\usepackage{xurl,hyperref}
\hypersetup{colorlinks=true,urlcolor=blue}
\setcounter{secnumdepth}{0}
\setlength{\parindent}{0pt}
\setlength{\parskip}{5pt}
\setlength{\emergencystretch}{3em}
\begin{document}
\section*{\texorpdfstring{Segre–Harbourne–Gimigliano–Hirschowitz Conjecture}{Segre–Harbourne–Gimigliano–Hirschowitz Conjecture}}\label{493}
\subsection{Definitions and mathematical statement}
For distinct general points $p_1,\ldots,p_r\in{\mathbb{P}}^2_{{\mathbb{C}}}$, let $\mathcal L(d;m_1,\ldots,m_r)$ be the projective space of degree-$d$ curves vanishing to order at least $m_i$ at $p_i$. Its virtual and expected dimensions are
\[
 v=\binom{d+2}{2}-1-\sum_i\binom{m_i+1}{2},
 \qquad e=\max\{-1,v\}.
\]
The system is special if its actual projective dimension exceeds $e$, with the empty system assigned dimension $-1$. On the blowup, a $(-1)$-curve is a smooth rational curve with self-intersection $-1$.

An exact reduced-form version of SHGH is that every Cremona-reduced system is nonspecial: for $m_1\ge m_2\ge\cdots\ge0$ and $d\ge m_1+m_2+m_3$, padding missing multiplicities by zero,
\[
 h^0(\mathcal L)=\max\left\{0,\binom{d+2}{2}
                                  -\sum_i\binom{m_i+1}{2}\right\}.
\]
Here $h^0(\mathcal L)$ means the dimension of the vector space before projectivizing. For each degree/multiplicity choice, general means points in an appropriate nonempty Zariski-open set. The geometric formulation uses $(-1)$-speciality, not an undefined fixed-locus condition for empty systems. The catalogue's URL is a paper by Brambilla--Postinghel, not the authors named in its citation; see \href{https://link.springer.com/article/10.1007/s40574-025-00468-5}{[E1]}.
\par\smallskip\textit{Formulation and concept references:} \href{https://link.springer.com/article/10.1007/s40574-025-00468-5}{[S1]}, \href{https://link.springer.com/article/10.1007/s40574-025-00468-5}{[E1]}.

\subsection{Short English statement}
Are unexpected dimensions of general plane-curve systems entirely explained by exceptional curves, so that every Cremona-reduced system has its expected dimension?
\subsection{Sources}
\begin{itemize}
\item[S1]Antonio Laface and Luca Ugaglia, Bollettino dell'Unione Matematica Italiana (2025).
\url{https://link.springer.com/article/10.1007/s40574-025-00468-5}.
\textit{Formulation source. Consulted 24 September 2026: Authorship is Brambilla–Postinghel, unlike the imported citation; Conjecture 3.1 is the SHGH statement.}
\item[E1]Maria Chiara Brambilla and Elisa Postinghel, Towards Good Postulation of Fat Points, One Step at a Time, Bollettino dell’Unione Matematica Italiana 18 (2025), 737–749, Sections 1 and 3.1, Conjecture 3.1.
\url{https://link.springer.com/article/10.1007/s40574-025-00468-5}.
\textit{Added primary source: Corrected bibliographic authorship and title for the URL already supplied by the catalogue; dimensions and the Cremona-reduced SHGH formulation. Consulted 24 September 2026.}
\end{itemize}
\end{document}
